\documentclass[11pt,a4paper]{article}

\usepackage[utf8]{inputenc}
\usepackage[T1]{fontenc}
\usepackage{amsmath,amssymb,amsthm}
\usepackage{booktabs}
\usepackage{graphicx}
\usepackage{microtype}
\usepackage[round]{natbib}
\usepackage[colorlinks=true, linkcolor=blue!60!black, citecolor=blue!60!black, urlcolor=blue!60!black]{hyperref}
\usepackage{xcolor}
\usepackage{listings}

\lstset{
  basicstyle=\ttfamily\small,
  breaklines=true,
  frame=single,
  framerule=0.2pt,
  language=Python,
  showstringspaces=false
}

\newcommand{\veropt}{\textsf{veropt}}
\newcommand{\real}{\mathbb{R}}
\newcommand{\E}{\mathbb{E}}

\title{Proxy-Informed Gaussian Process Priors for\\
Bayesian Optimisation of Very Expensive Objectives}

\author{James Avery\\
  \small Department of Computer Science, University of Copenhagen\\
  \small \texttt{avery@diku.dk}}

\date{June 2026 --- working draft}

\begin{document}

\maketitle

\begin{abstract}
Bayesian optimisation (BO) of very expensive objective functions --- e.g.\
coupled-cluster electronic structure energies that take months per evaluation
--- typically starts from an uninformative Gaussian process (GP) prior and
must spend a large fraction of its small evaluation budget rediscovering the
gross structure of the objective. In many such applications a fast surrogate
of the objective exists (harmonic approximations in microseconds,
semi-empirical methods in seconds) together with a \emph{known pointwise bound}
on its error: the true objective deviates from the proxy by at most a known
percentage, even though the locations of their optima may differ. We show how
to encode this knowledge exactly where it belongs in the BO machinery: the GP
prior mean is set to the proxy, and the prior covariance is scaled
heteroscedastically so that the prior standard deviation at every point is the
promised deviation band. The bound is treated as a hard upper limit on the
prior amplitude --- a single trainable scale factor may tighten the band as
data accumulates but is constrained never to exceed it. The construction
composes with arbitrary stationary correlation kernels, survives the data
renormalisation used in practical BO implementations, and is available in the
open-source package \veropt{}. On the six-dimensional Hartmann benchmark with
a proxy distorted by at most 1\%, the proxy-informed prior comes within
$0.3\%$ of the global optimum after a single Bayesian batch (12 evaluations)
on every tested seed, whereas a flat prior is still $19\%$ short of the
optimum when the budget of 24 evaluations is exhausted.
\end{abstract}

\section{Introduction}

Bayesian optimisation is the method of choice for optimising black-box
functions whose evaluations are so expensive that only a small number ---
often around one hundred --- can be afforded
\citep{shahriari2016, frazier2018}. A Gaussian process surrogate is fitted to
the evaluations made so far, and an acquisition function balances exploration
against exploitation to select the next batch of points.

The standard GP prior is maximally uninformative: a constant mean and a
stationary kernel with a single global amplitude. Every evaluation must then
contribute to learning both the gross structure of the objective \emph{and}
the fine structure near the optimum. For the most expensive objectives this is
wasteful, because the gross structure is often already known. Our motivating
example is molecular electronic structure: a CCSD(T)-level energy surface may
cost months of computation per point, while a harmonic approximation to the
same surface is available in microseconds and semi-empirical electronic
structure methods in seconds. Crucially, these cheap surrogates frequently
come with \emph{quantitative} error guarantees of the form ``the true energy
is within $1\%$ of the approximation, everywhere''. The minima of the cheap
and the expensive surface need not coincide --- otherwise the problem would be
solved --- but pointwise, the expensive function is trapped in a known band
around the cheap one.

This paper describes a simple, well-founded way to inject exactly this kind of
knowledge into BO: build the GP prior from the proxy. The prior mean is the
proxy itself, and the prior covariance is rescaled so that the prior standard
deviation at each point equals the promised band. The result is that the
surrogate starts out already knowing everything the proxy knows, and the
expensive evaluations are spent exclusively on learning the \emph{residual}
between truth and proxy --- a function whose amplitude is bounded by the
guarantee. Related ideas appear in multi-fidelity BO
\citep{kennedy2000, forrester2007}, where cheap and expensive evaluations are
modelled jointly, and in deep mean-function approaches; the variant presented
here is deliberately minimal: it requires no evaluations of the expensive
function to fit the relationship between fidelities, because that relationship
is supplied a priori as a hard bound.

\section{Method}

\subsection{Setting}

Let $f \colon X \subset \real^d \to \real$ be the expensive objective, to be
maximised under a budget of $N \sim 10^2$ evaluations. Let
$g \colon X \to \real$ be a fast proxy ($\mu$s--s per evaluation), and assume
the pointwise guarantee
\begin{equation}
  \label{eq:bound}
  \lvert f(x) - g(x) \rvert \;\le\; \varepsilon(x)
  \qquad \text{for all } x \in X,
\end{equation}
with $\varepsilon$ known. The two cases of practical interest are the
\emph{relative} bound $\varepsilon(x) = r\,\lvert g(x)\rvert$ (``$f$ is within
$1\%$ of $g$'', $r = 0.01$) and the \emph{absolute} bound
$\varepsilon(x) = \varepsilon_0$. Nothing is assumed about the locations of
the optima of $f$ and $g$.

\subsection{The proxy-informed prior}

We place on $f$ the Gaussian process prior
\begin{equation}
  \label{eq:prior}
  f \;\sim\; \mathcal{GP}\bigl(m, \, k\bigr),
  \qquad
  m(x) = g(x),
  \qquad
  k(x, x') = c^2\, \sigma(x)\, \sigma(x')\, \rho(x, x'),
\end{equation}
where $\rho$ is any correlation kernel ($\rho(x,x) = 1$; e.g.\ Mat\'ern-5/2,
rational quadratic, or sums thereof) and
\begin{equation}
  \label{eq:sigma}
  \sigma(x) \;=\; \frac{\varepsilon(x)}{\kappa},
\end{equation}
with $\kappa > 0$ a confidence parameter discussed below. The kernel in
\eqref{eq:prior} is the standard non-stationary scaling of a correlation
kernel by an outer product of positive functions, and is therefore a valid
(positive semi-definite) covariance.

Equation~\eqref{eq:prior} is most easily read through the residual
$h = f - g$: the prior states $h \sim \mathcal{GP}(0, k)$ with pointwise
standard deviation $c\,\sigma(x)$, i.e.\ the model believes the truth deviates
from the proxy by a smooth function whose typical size at $x$ is a known
fraction of the band $\varepsilon(x)$. With $c = 1$ and $\kappa = 2$ the prior
places ${\sim}95\%$ of its mass inside the band \eqref{eq:bound} at every
point; $\kappa = 3$ gives ${\sim}99.7\%$. The bound is thus encoded
\emph{softly} --- a GP cannot represent a hard truncation --- and $\kappa$
expresses how literally the guarantee is taken. We default to $\kappa = 2$.

\subsection{The bound as a hard cap on a trained amplitude}

The guarantee \eqref{eq:bound} is an \emph{upper} bound: the proxy may well be
better than its worst case. We therefore make the global amplitude factor $c$
in \eqref{eq:prior} a hyperparameter, trained alongside the kernel
lengthscales by type-II maximum likelihood, but constrained to
\begin{equation}
  c \in (c_{\min}, \, 1],
\end{equation}
implemented by a sigmoidal reparametrisation (a \texttt{gpytorch}
\citep{gardner2018} interval constraint). The cap at $1$ guarantees that no
amount of hyperparameter optimisation can inflate the prior beyond the
promised band, while the marginal likelihood is free to \emph{tighten} the
band as the data reveal that the proxy is more accurate than worst-case ---
which immediately sharpens the posterior everywhere, not only near
evaluations. The floor $c_{\min}$ (default $0.1$) prevents overconfident
collapse from the handful of points available early on. Setting $c$ constant
at $1$ recovers the literal worst-case prior.

Two practical guards are worth stating. First, under a relative bound
$\sigma(x) \to 0$ wherever $g(x) \to 0$, which pins the GP to the proxy and
can destabilise Cholesky factorisations; a small floor
$\varepsilon(x) \gets \max(\varepsilon(x), \varepsilon_{\mathrm{floor}})$
restores numerical safety. Second, since $\rho(x,x) = 1$ exactly, the prior
variance at any point is $c^2 \sigma(x)^2$ by construction --- there is no
additional global outputscale to interfere with the calibration.

\subsection{Interaction with data normalisation}

Practical BO implementations whiten their data: inputs and observed objective
values are affinely normalised (zero mean, unit variance per dimension), and
the normalisers are refitted as new evaluations arrive. The proxy, by
contrast, is defined in physical units. The prior mean must therefore be
evaluated as
\begin{equation}
  \tilde m(\tilde x) \;=\; T_y\bigl(g(T_x^{-1}(\tilde x))\bigr),
\end{equation}
where $T_x, T_y$ are the \emph{current} input and output normalisations and
tildes denote whitened quantities. The band, being the standard deviation of a
\emph{deviation} rather than a value, transforms through the scale of $T_y$
only:
\begin{equation}
  \tilde\sigma(\tilde x) \;=\; T_y\bigl(\sigma(x)\bigr) - T_y(0)
  \;=\; \sigma(x) / s_y,
\end{equation}
with $s_y$ the output scale --- applying the full affine map here would be
incorrect, silently shifting the band by the normalisation offset. Because the
normalisers change between refits, the implementation binds them to the mean
and kernel by late-bound closures that read the current normaliser at call
time.

\subsection{Cost}

The proxy is evaluated at every posterior query, which during acquisition
optimisation (here: dual annealing with sequential proximity-punished batch
selection) means thousands of calls per suggested batch. Microsecond- to
millisecond-scale proxies are negligible; second-scale proxies (semi-empirical
methods) dominate the suggestion step and benefit from memoisation, which we
leave as future work. Against an objective costing days to months per
evaluation, even the un-memoised cost is irrelevant.

\section{Implementation}

The method is implemented in the open-source Bayesian optimisation package
\veropt{} \citep{veropt}, which was built for tuning the VEROS ocean simulator
\citep{hafner2018} and targets exactly the expensive-objective regime
considered here. The user subclasses a savable \texttt{ProxyMeanFunction}
(physical units in, physical units out) and passes it through the model
configuration:

\begin{lstlisting}
optimiser = bayesian_optimiser(
    n_initial_points=8,
    n_bayesian_points=32,
    n_evaluations_per_step=4,
    objective=expensive_objective,
    model={
        'kernels': 'matern',
        'proxy_prior': my_proxy,
        'proxy_prior_settings': {'bound_value': 0.01},
    },
    n_points_before_fitting=4,
)
\end{lstlisting}

Internally, a \texttt{ProxyMean} module supplies the mean of \eqref{eq:prior},
and a \texttt{Proxy\-Scaled\-Kernel} wraps the chosen correlation kernel with
the $c^2\sigma(x)\sigma(x')$ factor; the construction composes with every kernel
offered by the package and with its per-step renormalisation, batch
acquisition, and state serialisation. In the multi-objective case a proxy may
be supplied per objective, with uninformed objectives left at the default
prior.

\section{Experiments}

We use the six-dimensional Hartmann function (maximised, optimum
$f^* = 3.3224$) as the ``expensive'' objective and construct a proxy by
multiplicative harmonic distortion,
\begin{equation}
  g(x) \;=\; f(x)\,\bigl(1 + 0.01 \sin(6\pi x_1)\bigr),
\end{equation}
so that the relative bound \eqref{eq:bound} with $r = 0.01$ holds exactly,
while the distortion shifts the local geometry of the optimum. Both methods
use a Mat\'ern-5/2 kernel with ARD, batches of $4$, $8$ initial random points
followed by $16$ Bayesian-phase points (a total budget of $24$ evaluations),
and \veropt{}'s defaults otherwise. The proxy-informed runs use $\kappa = 2$,
a trainable amplitude $c \in (0.1, 1]$, and an absolute band floor of
$\varepsilon_{\mathrm{floor}} = 0.01$. Each configuration is repeated over
five random seeds.

\begin{table}[t]
  \centering
  \caption{Best objective value found (mean $\pm$ standard deviation over five
  seeds) after $n$ evaluations of the Hartmann-6 objective
  ($f^{*} = 3.3224$). Both methods share seeds, budget, kernel, and
  acquisition machinery; they differ only in the prior.}
  \label{tab:results}
  \begin{tabular}{lcc}
    \toprule
    $n$ & Flat prior & Proxy-informed prior \\
    \midrule
    $8$ (random)  & $1.117 \pm 0.644$ & $1.117 \pm 0.644$ \\
    $12$          & $1.806 \pm 0.491$ & $3.3152 \pm 0.0010$ \\
    $16$          & $2.052 \pm 0.155$ & $3.3178 \pm 0.0041$ \\
    $20$          & $2.571 \pm 0.186$ & $3.3178 \pm 0.0041$ \\
    $24$          & $2.679 \pm 0.190$ & $3.3178 \pm 0.0041$ \\
    \bottomrule
  \end{tabular}
\end{table}

Table~\ref{tab:results} reports the best objective value found as a function
of the number of evaluations $n$, over five shared seeds (the $8$ initial
random points are identical between the two arms). The effect of the prior is
immediate: after a \emph{single} Bayesian batch ($n = 12$), every
proxy-informed run sits within $0.008$ of the global optimum --- a relative
optimality gap below $0.3\%$, i.e.\ at the scale of the proxy's own error
band --- and the across-seed spread collapses from $\pm 0.49$ to
$\pm 0.001$. The flat prior, by contrast, is still climbing at the end of the
budget and terminates $0.64$ below the optimum on average (a $19\%$ gap) with
fifty times the seed-to-seed variability. Reading the comparison in terms of
evaluations rather than attained value: the proxy-informed optimiser needs
$12$ evaluations to surpass the level the flat prior has not yet reached after
$24$. For an objective costing months per evaluation, this is the difference
between one and several years of computation. We note the benchmark is
deliberately favourable --- the proxy is everywhere within $1\%$ of the truth
and the distortion is smooth --- but this is precisely the regime the method
is designed for, and the guarantee that makes it favourable is one the user
must certify, not one the method assumes silently.

\section{Limitations and future work}

\paragraph{Soft bound.} The hard guarantee \eqref{eq:bound} is encoded as a
$\kappa$-sigma Gaussian band; at $\kappa = 2$, roughly $5\%$ of prior mass
violates the bound pointwise. Truncated or constrained GPs would encode the
bound exactly at considerably greater inferential cost.

\paragraph{Calibration of the residual smoothness.} The bound fixes the
\emph{amplitude} of the residual but not its \emph{lengthscale}; the latter is
learned from data as usual. A proxy whose error oscillates faster than the
learned lengthscale will be modelled overconfidently between evaluations.

\paragraph{Zero-data suggestion.} With an informative prior the surrogate is
meaningful before any expensive evaluation, so in principle the initial random
phase could be replaced by acquisition over the prior alone. The current
implementation still requires one batch of evaluations before the model is
activated (its normalisers need at least two points); removing this
restriction is mechanical but not yet done.

\paragraph{Gradients and memoisation.} The proxy is currently evaluated
without gradients --- harmless under derivative-free acquisition optimisation,
but a gradient-based acquisition optimiser would require differentiating
through the prior mean. Memoisation of proxy calls would make second-scale
proxies practical.

\section{Conclusion}

When an expensive objective comes with a cheap surrogate and a quantitative
pointwise error guarantee --- the everyday situation in computational
chemistry and many other simulation sciences --- that guarantee can be
installed directly into the Gaussian process prior of a Bayesian optimiser:
proxy as prior mean, error band as (heteroscedastic) prior standard deviation,
bound as a hard cap on a trained amplitude. The construction is exact in its
bookkeeping (valid PSD covariance, correct behaviour under data whitening),
composes with standard kernels and acquisition machinery, and in benchmark
experiments converts a months-per-point optimisation that would otherwise
spend its budget rediscovering known structure into one that starts at the
proxy and spends every evaluation on the residual.

\section*{Acknowledgements}

The proxy-informed prior is implemented in \veropt{}, developed by Aster
Stoustrup and collaborators; this work builds directly on that foundation.

\begin{thebibliography}{9}

\bibitem[Forrester et~al.(2007)]{forrester2007}
A.~I.~J. Forrester, A.~S\'obester, and A.~J. Keane.
\newblock Multi-fidelity optimization via surrogate modelling.
\newblock \emph{Proceedings of the Royal Society A}, 463:3251--3269, 2007.

\bibitem[Frazier(2018)]{frazier2018}
P.~I. Frazier.
\newblock A tutorial on Bayesian optimization.
\newblock \emph{arXiv:1807.02811}, 2018.

\bibitem[Gardner et~al.(2018)]{gardner2018}
J.~R. Gardner, G.~Pleiss, D.~Bindel, K.~Q. Weinberger, and A.~G. Wilson.
\newblock GPyTorch: Blackbox matrix-matrix Gaussian process inference with GPU
  acceleration.
\newblock In \emph{Advances in Neural Information Processing Systems}, 2018.

\bibitem[H{\"a}fner et~al.(2018)]{hafner2018}
D.~H{\"a}fner, R.~L. Jacobsen, C.~Eden, M.~R.~B. Kristensen, M.~Jochum,
  R.~Nuterman, and B.~Vinter.
\newblock Veros v0.1 -- a fast and versatile ocean simulator in pure Python.
\newblock \emph{Geoscientific Model Development}, 11:3299--3312, 2018.

\bibitem[Kennedy and O'Hagan(2000)]{kennedy2000}
M.~C. Kennedy and A.~O'Hagan.
\newblock Predicting the output from a complex computer code when fast
  approximations are available.
\newblock \emph{Biometrika}, 87(1):1--13, 2000.

\bibitem[Shahriari et~al.(2016)]{shahriari2016}
B.~Shahriari, K.~Swersky, Z.~Wang, R.~P. Adams, and N.~de~Freitas.
\newblock Taking the human out of the loop: A review of Bayesian optimization.
\newblock \emph{Proceedings of the IEEE}, 104(1):148--175, 2016.

\bibitem[Stoustrup et~al.(2026)]{veropt}
A.~Stoustrup et~al.
\newblock \veropt{}: the versatile optimiser.
\newblock \url{https://github.com/aster-stoustrup/veropt}, 2026.

\end{thebibliography}

\end{document}
